\scanpage{182}
% manuscript position: top
\mspage{98}
% manuscript position: 1–3
stable par $\underline{\lim}$ de cardinal $\leq\pi$ et b) $M$ est stable par $\varinjlim$ \inserted{sur des ens ord $I$} grands devant $\pi$. Alors $M$ est stable par petites $\varinjlim$.
% manuscript position: left/1–9
% margin note written sideways
\margin{\uline{Cor. 4.4.14.1} $M$ étant comme dans 4.4.14, et $f:M\longrightarrow\unsure{\mathcal E|E}$ un foncteur qui commute aux \inserted{petites} $\varinjlim$ de cardinal $\leq\pi$ et aux $\varinjlim$ \deleted{\unsure{qq}} \unsure{grandes devant} $\pi$, alors $f$ commute aux petites $\varinjlim$.}

% manuscript position: 4–7
\deleted{\uline{NB} On sait que ces hypothèses sont vérifiées pour $\operatorname{Ind}_\pi C$, sous les conditions 4.4.13, grâce à 4.4.13 b) déjà établi et \deleted{\unsure{au}} lemme 4.4.8.1}

% manuscript position: 8–9
En effet, si $\varphi:J\longrightarrow M$ est donné, on en \unsure{conclut|construit}
% manuscript position: 10
\[
\psi:I=I(J)\longrightarrow M
\]
% manuscript position: 11
par
% manuscript position: 12
\[
\psi(J')=\varinjlim_{J'}\varphi\mid J',
\]
% manuscript position: 13–18
qui existe par hypothèse sur $M$. Il est alors immédiat que $\varinjlim_J\varphi$ existe sss $\varinjlim_I\psi$ existe, or cette existence est assurée du fait que $I$ est grand devant $\pi$, et l'hypothèse sur $M$.

% manuscript position: 19–22
Pour achever de prouver le \unsure{th.} 4.4.13, \deleted{\unsure{on}} est donc \unsure{amené|ramené} au cas particulier de limites inductives $\varinjlim_J\varphi(j)$, qd $\varphi:J\deleted{\to}\longrightarrow\operatorname{Ind}'_\pi C$ est donné, avec $J$ \unsure{ens|une}
% manuscript position: foot
% archive sheet number: 181
